\scanpage{126}\mspage{74}
% paper clip at upper right
\[
\handwritten{H^1(\mathscr E_{\varphi},O_{\mathscr E_{\varphi}}[nE]_{\mathscr E_{\varphi}})}=\handwritten{0}
\]
we obtain an exact sequence
\[
\handwritten{H^{\unsure{0|o}}(V,O_V[nE])\xrightarrow{r^*}H^0(\mathscr E_{\varphi},O_{\mathscr E_{\varphi}}[nE]_{\mathscr E_{\varphi}})\to0}\tag*{\handwritten{(19)}}
\]

Let us assume that there would exist a point $P\in\mathscr E_{\varphi}$
such that in $P$ any cross-section $\Phi\in H^0(V,O_V[nE])$
satisfies the condition: if $P\in U^{(i)}$, then $\Phi^{(i)}(P)=0$.
Since any cross-section $\Phi'\in H^0(\handwritten{\mathscr E_{\varphi}},\handwritten{O}_{\deleted{\unsure{\varphi|o}}\unsure{\mathscr E_{\varphi}|\varphi}}[nE]_{\mathscr E_{\varphi}})$ can be
presented in the form $\Phi'=r^*\Phi$ or $\Phi^{\prime(i)}=r\,\Phi^{(i)}$ and
$r$ in the stalk over $P$ is a homomorphism of restriction
of functions; then $\Phi^{\prime(i)}(P)=0$. But this contradicts
the fact that, at $n\ge n(\mathscr E_{\varphi})$, the elements of the group
$H^0(\mathscr E_{\varphi},O_{\mathscr E_{\varphi}}[nE]_{\mathscr E_{\varphi}})$ generate the stalk of the sheaf
$O_{\mathscr E_{\varphi}}[nE]_{\mathscr E_{\varphi}}$ over any point $P\in\mathscr E_{\varphi}$. The lemma is proved. $\blacksquare$

Let us come back to the proof of theorem 2. Put
$E'=n_{\mathscr E}E$, where $n_{\mathscr E}$ is the number to which lemma 3
refers. Let $\Phi_{\unsure{0|o}},\ldots,\Phi_T$ be the basis (not necessarily
of linearly independent vectors) of the group
$H^0(V,O_V[nE'])$.

Consider a mapping $f_n\colon V\to P^T$ ($P^T$ is a
$T$-dimensional projective space), operating according to
the rule: if $P\in V$, $P\in U^{(i)}$, then
% archive sheet number: 125
